\section{Local Answer-Cell Generator Theorem}

The previous sections introduced Lift \& Twist as a coordinate-straightening mechanism.  The Project 22 audit sequence now gives a stronger local result: within the four-state reduced target, the answer-cell data are generated from two bits.

Let \(b\in\{0,1\}\) be the shell bit, with \(b=0\) for the ordinary shell and \(b=1\) for the branch shell.  Let \(r\in\{0,1\}\) be the rank bit.  The four local states are
\[
    O0=(0,0),\quad
    O1=(0,1),\quad
    B0=(1,0),\quad
    B1=(1,1).
\]

\begin{theorem}[Local Lift \& Twist answer-cell theorem]
For the four-state local answer-cell target, the state bits \(b,r\) generate the C row, anchor column, selected candidate, C key, anchor key, C path, C value set, anchor residue set, overlap markers, and closed anchor node path.  The generated table agrees exactly with the Project 22 target table.
\end{theorem}

\paragraph{Header laws.}
The visible placement and key data are generated by
\[
    c_{\mathrm{row}} = 2b+r,
\]
\[
    a_{\mathrm{col}} = 1+r+b(2r-1),
\]
\[
    s = 4c_{\mathrm{row}} + a_{\mathrm{col}},
\]
\[
    c_{\mathrm{key}} \equiv 11+6b+2r \pmod{15},
\]
\[
    a_{\mathrm{key}} \equiv 4b+12r+12br \pmod{15}.
\]

\paragraph{C-side path law.}
Let
\[
    d_1 \equiv 6-3b-3r+br \pmod{15},
\]
\[
    d_2 \equiv (1-b)(3+9r)+13b \pmod{15}.
\]
Then the C path is
\[
    [c_{\mathrm{key}},\ c_{\mathrm{key}}+d_1,\ c_{\mathrm{key}}+d_2,\ c_{\mathrm{key}}]\pmod{15}.
\]

\paragraph{Anchor-side residue law.}
The anchor key is \(a_{\mathrm{key}}\).  Ordinary states use offsets
\[
    \{5,7,12\}\cup\{1+3(1-r),\ 2+6(1-r),\ 3+6(1-r)\},
\]
while branch states use offsets
\[
    \{0,4\}\cup\{11-6r,\ 13-4r,\ 14-4r\},
\]
with an additional terminal offset \(12\) when \(r=1\).  Adding the appropriate offsets to \(a_{\mathrm{key}}\) modulo 15 gives the anchor residue set.

\paragraph{Closed anchor path law.}
The ordered anchor paths are generated by ordered offsets and shell-dependent lift masks.  Ordinary states use the open lift mask
\[
    [1,1,0,0,0,0],
\]
and branch states use the open lift mask
\[
    [0,0,0,1,0,1].
\]
The open path is then closed by returning to the first node-pair.

\paragraph{Generated target.}
\[
\begin{array}{c|c|c|c|c|c|c|c|c|c}
\mathrm{state} & b & r & s & c_{\mathrm{key}} & a_{\mathrm{key}} & C\ \mathrm{path} & C\ \mathrm{values} & A\ \mathrm{residues} & \mathrm{overlap}\\
\hline
O0 & 0 & 0 & 1 & 11 & 0 & \texttt{11 2 14 11} & \texttt{2 11 14} & \texttt{4 5 7 8 9 12} & \texttt{none} \\
O1 & 0 & 1 & 6 & 13 & 12 & \texttt{13 1 10 13} & \texttt{1 10 13} & \texttt{0 2 4 9 13 14} & \texttt{13} \\
B0 & 1 & 0 & 8 & 2 & 4 & \texttt{2 5 0 2} & \texttt{0 2 5} & \texttt{0 2 3 4 8} & \texttt{0 2} \\
B1 & 1 & 1 & 15 & 4 & 13 & \texttt{4 5 2 4} & \texttt{2 4 5} & \texttt{2 3 7 8 10 13} & \texttt{2} \\
\end{array}
\]

Artifacts 007--012 verify these laws independently and then consolidate them.  Artifact 011 verifies that the generated rows match the local answer-cell target.  Artifact 012 verifies the closed anchor node path geometry.

\paragraph{Boundary.}
This theorem is local to the four-state reduced target.  It does not yet derive the full reduced 16-candidate universe from native G60 provenance.  It does not derive a full G60-native generator, and it does not close Gap A.
